\chapter{Yu Deng}\label{ch2026yudeng:yu-deng-fields-medal-2026}

\begin{center}\rule{0.5\linewidth}{0.5pt}\end{center}

\textbf{Prepared:} July 2026 \textbf{Subject:} Yu Deng (b. 1990, Shenzhen, Guangdong, China) \textbf{Position:} Professor of Mathematics, University of Chicago \textbf{Doctoral Advisor:} Alexandru D. Ionescu (Princeton University) \textbf{Fields Medal Citation:} \emph{``For his work in partial differential equations, including the rigorous derivation of the Boltzmann equation from hard-sphere dynamics for rarefied gases, the derivation of wave kinetic equations from nonlinear dispersive systems, and probabilistic approaches to nonlinear Schr\"odinger dynamics.''}

\begin{center}\rule{0.5\linewidth}{0.5pt}\end{center}

\begin{prereqs} Partial differential equations, statistical mechanics, Newtonian mechanics, functional analysis, probability theory. \end{prereqs}

\begin{objectives} Explain Hilbert's sixth problem and its modern formulation; describe the Boltzmann equation and its role in kinetic theory; understand the strategy for deriving the Boltzmann equation from particle dynamics \cite{denghanima2024}. \end{objectives}

\section{Problem History}

In 1900, David Hilbert presented his famous list of 23 unsolved problems. The sixth problem called for the axiomatization of physics: ``To treat in the same manner, by means of axioms, those physical sciences in which already today mathematics plays an important part; in the first rank are the theory of probabilities and mechanics.'' The modern interpretation asks whether the macroscopic kinetic equations governing gases---such as the Boltzmann equation---can be derived from the microscopic Newtonian dynamics of a large number of particles.

The Boltzmann equation, introduced by Ludwig Boltzmann in 1872, describes the statistical behavior of a gas at the kinetic level:
\[
\partial_t f + v \cdot \nabla_x f = Q(f,f),
\]
where $f(t,x,v)$ is the distribution of particles at position $x$ and velocity $v$, and $Q(f,f)$ is the collision operator encoding binary elastic collisions. The equation rests on the \emph{molecular chaos assumption} (Stosszahlansatz): colliding particles are assumed to have uncorrelated velocities before collision.

Lanford proved in 1975 a rigorous derivation of the Boltzmann equation from Newtonian dynamics for hard spheres on a short time interval. Extending the derivation to longer times had been a major open problem.

The Boltzmann equation is a nonlinear integro-differential equation that describes the evolution of the distribution function $f(t,x,v)$. The collision operator $Q(f,f)$ is defined by
\[
Q(f,f)(v) = \int_{\mathbb{R}^3} \int_{S^2} (f(v') f(w') - f(v) f(w)) |v-w| \, d\sigma \, dw,
\]
where $v'$ and $w'$ are the post-collision velocities determined by the conservation of momentum and energy. The equation conserves mass, momentum, and energy, and the $H$-theorem states that the Boltzmann entropy $H(f) = \int f \log f \, dx \, dv$ is non-increasing along solutions.

The fundamental difficulty in deriving the Boltzmann equation from Newtonian dynamics is the \emph{propagation of molecular chaos}. The molecular chaos assumption states that before a collision, the velocities of the two colliding particles are uncorrelated: $\langle f_2(x_1, v_1, x_2, v_2) \rangle = f(x_1, v_1) f(x_2, v_2)$, where $f_2$ is the two-particle distribution function. This assumption is justified if the particles are well-separated and have not interacted recently, but after many collisions, the correlations between particles can build up and threaten the validity of the assumption.

The mathematical challenge is to show that these correlations remain controlled as time evolves. Deng, Hani, and Ma developed a long-time cumulant approach that applies throughout the lifespan of the corresponding Boltzmann solution.

Wave turbulence theory provides a parallel story. It predicts that energy injected at one scale in a system of interacting waves will cascade to other scales through nonlinear interactions. The \emph{wave kinetic equation}---the analogue of the Boltzmann equation for wave systems---describes this energy cascade. Deriving the wave kinetic equation from the underlying nonlinear wave equations with rigorous error bounds was another fundamental problem in mathematical physics.

The difficulty in both cases is that repeated interactions create correlations between particles (or waves), threatening the molecular chaos assumption. A single collision between two particles is simple to analyze, but after many collisions, the state of the system becomes highly correlated. Controlling these correlations---showing they remain small on average---is the central mathematical challenge.

\section{How Was It Solved}

In their 2024 paper, Deng, Hani, and Ma prove a long-time derivation for hard spheres in the three-dimensional Boltzmann--Grad scaling, with $N\epsilon^2$ fixed. More precisely, the result holds on any time interval on which the limiting Boltzmann solution exists; it is not an unconditional statement for all positive times \cite{denghanima2024}.

Their argument propagates a cumulant ansatz that retains information about the particles' collision histories. The estimates reduce to combinatorial properties of collision-history diagrams, called collision-history molecules in the paper. A cutting algorithm controls these diagrams and is a central new ingredient. This description does not assert a universal factorial decay law for $k$-fold collisions or uniform-in-time error estimates independent of the solution's lifespan.

In related work, Deng and Hani developed methods for deriving wave kinetic equations from nonlinear dispersive equations. This work and the hard-sphere result share an emphasis on controlling long interaction histories, but the proofs use different models and estimates.

The work on wave turbulence preceded and motivated the attack on Hilbert's sixth problem. The techniques for controlling repeated interactions in wave systems were adapted to the particle setting, where the analogous problem---controlling repeated collisions---is even more challenging because particles are discrete objects rather than continuous fields.

The wave kinetic equation describes the evolution of the wave energy spectrum $n(t,k)$ in a system of weakly interacting waves:
\[
\partial_t n = Q(n,n),
\]
where the collision operator $Q$ encodes the nonlinear wave interactions. Deriving this equation from an underlying nonlinear wave equation requires controlling correlations among interacting modes.

The wave and particle settings have different interaction laws and require distinct analytic estimates, even though both problems involve the accumulation of correlations over time.

\section{Impact of This Discovery}

The result extends Lanford's short-time theorem to intervals of arbitrary length, subject to the existence of the limiting Boltzmann solution on those intervals. It is a major advance in the rigorous derivation of kinetic equations from microscopic dynamics.

The derivation validates the molecular chaos assumption---Boltzmann's key postulate---by showing it emerges naturally from Newtonian dynamics in the appropriate limit. This closes a foundational gap that had persisted since the earliest days of statistical mechanics.

The hard-sphere theorem leaves open how far its methods can extend to other interaction laws and kinetic limits. Such extensions require additional arguments and are active research questions.

Deng's work on wave turbulence has contributed to the mathematical understanding of energy cascades in physical systems. The rigorous derivation of wave kinetic equations has potential applications in climate science (understanding energy transfer in atmospheric waves), plasma physics (wave-particle interactions), and quantum field theory (thermalization of quantum systems).

The proof demonstrates the power of modern PDE techniques---including dispersive estimates, harmonic analysis, probability theory, and kinetic theory---in addressing fundamental questions about the mathematical foundations of physics.

The \textbf{BBGKY hierarchy} (Bogoliubov--Born--Green--Kirkwood--Yvon) provides a systematic expansion of the many-particle dynamics in terms of reduced distribution functions. The hierarchy couples the $k$-particle distribution function to the $(k+1)$-particle distribution function, and truncating at the first level gives the Boltzmann equation under the molecular chaos assumption. Deng's work provides a rigorous justification for this truncation by showing that the correlations between particles remain small in the Boltzmann--Grad limit.

\section{Modern Views and Improvements}

The derivation of the Boltzmann equation has opened numerous research directions. Extensions to softer interaction potentials (Lennard-Jones, Coulomb) would cover physically realistic gases. The derivation for quantum systems---replacing Newtonian dynamics with the Schr\"{o}dinger equation---is a major open problem.

The \textbf{Boltzmann--Grad limit} for hard spheres in three dimensions takes $N\to\infty$ and particle diameter $\epsilon\to0$ while keeping $N\epsilon^2$ fixed. Deng, Hani, and Ma establish the derivation on time intervals within the lifespan of the limiting Boltzmann solution, controlling the correlations through their cumulant and collision-history analysis \cite{denghanima2024}.

The long-time result improves on Lanford's short-time interval. Its scope is governed by the lifespan of the limiting solution; the paper should be consulted for the theorem's precise quantitative hypotheses and error estimates.

Related work studies long-time derivations of wave kinetic equations from nonlinear dispersive systems. These results use methods tailored to the wave setting.

Future directions include the derivation of hydrodynamic equations (Navier--Stokes, Euler) from kinetic theory, the study of fluctuations and large deviations in the Boltzmann--Grad limit, and the application of these techniques to quantum many-body physics.

The wave and particle results illustrate how related questions about long-time correlations arise in different physical models, while requiring distinct mathematical frameworks.

Deng and Hong Wang were classmates at Peking University, making 2026 the first time two former classmates have simultaneously won the Fields Medal. Their shared mathematical interests in harmonic analysis and PDE theory, developed during their undergraduate years, led to parallel breakthroughs in seemingly different areas---Wang in geometric analysis and Deng in kinetic theory. Their collaboration and mutual influence exemplify the interconnectedness of modern mathematics.

The rigorous derivation of the Boltzmann equation has implications for the foundations of statistical mechanics. The second law of thermodynamics states that the entropy of an isolated system is non-increasing, and this law is derived from the $H$-theorem of the Boltzmann equation. Deng's rigorous derivation provides a mathematical foundation for the second law by showing that it emerges from the Newtonian dynamics of particles. This closes a foundational gap that had persisted since the earliest days of statistical mechanics, and it provides a rigorous justification for the use of kinetic equations in physics.

The connection between kinetic theory and hydrodynamic limits is an active direction. The Euler and Navier--Stokes equations can be derived from the Boltzmann equation in the hydrodynamic limit, where the mean free path of particles goes to zero. Deng's rigorous derivation of the Boltzmann equation provides a starting point for the derivation of hydrodynamic equations from first principles, which would complete the chain from Newtonian mechanics to macroscopic fluid dynamics.

The study of the Boltzmann equation for quantum systems is another major open problem. In quantum mechanics, the dynamics of particles is governed by the Schrodinger equation rather than Newton's laws, and the derivation of the quantum Boltzmann equation from the Schrodinger equation requires new techniques. The quantum Boltzmann equation describes the evolution of the Wigner function, which is the quantum analogue of the classical distribution function. The derivation of the quantum Boltzmann equation from first principles is a major challenge that connects Deng's work on classical kinetic theory to quantum mechanics.

Questions about extending long-time derivations to other interaction laws, studying fluctuations around kinetic limits, and connecting kinetic to hydrodynamic equations remain open research directions. The cited hard-sphere theorem does not establish results for Vlasov or quantum systems, or for large deviations.

The wave kinetic equation derived by Deng and Hani has applications to the study of wave turbulence in physical systems. Wave turbulence describes the statistical behavior of systems of weakly interacting waves, such as ocean waves, plasma waves, and optical waves. The wave kinetic equation describes the evolution of the wave energy spectrum, and the rigorous derivation of this equation from the underlying nonlinear wave equation provides a mathematical foundation for the theory of wave turbulence.

The study of the Boltzmann equation with boundary conditions is another active direction. The interaction of a gas with a solid boundary is described by boundary conditions on the Boltzmann equation, and the derivation of these boundary conditions from the microscopic dynamics is a major open problem. The development of rigorous methods for studying the Boltzmann equation with boundary conditions, including the derivation of the Knudsen boundary condition and the accommodation coefficient, is an active area that builds on Deng's foundational work on the derivation of the Boltzmann equation.

\begin{center}\rule{0.5\linewidth}{0.5pt}\end{center}

\section*{Summary}
\addcontentsline{toc}{section}{Summary}

Deng, with Hani and Ma, rigorously derived the Boltzmann equation from hard-sphere dynamics on time intervals within the lifespan of the limiting solution, a major step toward Hilbert's sixth problem. His work also includes derivations of wave kinetic equations from nonlinear dispersive systems.
